\section{Research Notes and Selected Sources}

\subsection{Methodological Framing}

This chapter combines empirical research in organizational economics and entrepreneurship with original heuristics designed for founder self-diagnosis. The inventory tables, evidence levels E0 to E5, 6-dimension founder-market fit rubric, and 30-day discovery protocol are decision aids, not calibrated psychological psychometrics. They are designed to organize evidence, expose ignorance, and structure bounded tests.

\subsection{Selected Empirical References}

\begin{enumerate}[label=\textbf{[\arabic*]},leftmargin=11pt,itemsep=2.5pt,topsep=1.5pt,parsep=0pt]
  \item \textbf{Saras D. Sarasvathy (2001).} ``Causation and Effectuation: Toward a Theoretical Shift from Economic Inevitability to Entrepreneurial Contingency.'' \textit{Academy of Management Review} 26, no. 2: 243--263. Formalizes the bird-in-hand principle: starting from identity, knowledge, and network ties.
  \item \textbf{John Kitching and Julia Rouse (2020).} ``Contesting effectuation theory: Why it does not explain new venture creation.'' \textit{International Small Business Journal} 38, no. 6: 515--536. A rigorous critique emphasizing structural constraints and market structures that effectuation understates.
  \item \textbf{Scott Shane (2000).} ``Prior Knowledge and the Discovery of Entrepreneurial Opportunities.'' \textit{Organization Science} 11, no. 4: 448--469. Shows how prior idiosyncratic knowledge determines which applications entrepreneurs see in a single technological invention.
  \item \textbf{Meta-Analysis on Prior Experience (2021).} ``Does prior experience matter? A meta-analysis of the relationship between prior experience of entrepreneurs and firm performance.'' \textit{Journal of Small Business Management} 61, no. 4: 1851--1886. Synthesizes 85 studies; finds modest correlation (roughly 0.095) with significant moderator effects.
  \item \textbf{Karthik Rajkumar et al. (2022).} ``A causal test of the strength of weak ties.'' \textit{Science} 377, no. 6612: 1304--1310. A massive algorithmic experiment on LinkedIn demonstrating causal effects of tie strength on job mobility.
  \item \textbf{Paul Graham (2012).} ``How to Get Startup Ideas.'' Available online via \href{https://www.paulgraham.com/startupideas.html}{\color{BayesTeal}paulgraham.com}. Practitioner insight on living at the leading edge of a domain.
  \item \textbf{Pierre Azoulay, Benjamin F. Jones, J. Daniel Kim, and Javier Miranda (2020).} ``Age and High-Growth Entrepreneurship.'' \textit{American Economic Review: Insights} 2, no. 1: 65--82. Demonstrates with U.S. administrative data that peak founder performance occurs around age 45, driven by accumulated industry knowledge.
  \item \textbf{Saras D. Sarasvathy (2006).} ``The Bird-in-Hand Principle: Who I Am, What I Know, and Whom I Know.'' UVA Darden Technical Note. Practical operational guide for starting from existing capabilities.
\end{enumerate}

\subsection{What This Chapter Does Not Establish}

This chapter does not prove that a specific founder will triumph; that a target market has paying demand; that a given credential creates pricing power; or that any initial skill will become an unassailable competitive moat. Those are empirical hypotheses that demand customer discovery, unit economics, and competitive testing.

\subsection{The 1-Page Deliverable}

\textbf{When you finish this chapter, you should be able to state, with verifiable receipts: what you can do, what you know, who will respond, what you can afford, when you can work, what you cannot use, which problems you are positioned to investigate, and what experiment you will run next.}
